% =====================================================================
% VERSION ENSAE, REDUCTION SFE : conclusion.
% Phase 3 : copie de chapitres/14_conclusion.tex. Phase 4 (1er octobre 2026) : reecrite selon la
% regle n°2 de la mission et ramenee a deux pages, en suivant les consignes MS (situer le point
% d'arrivee, ce qui a pu etre realise, ce qui n'a pu l'etre, la suite possible). Rapport :
% reduction/chapitres/conclusion.md, qui liste ce qui est parti et ou.
% Reprise le 1er octobre 2026 tard le soir (« une ecriture plus humaine ») : registre cale sur
% des memoires d'actuariat publies avant 2022 (style/registre_humain.md), longueur inchangee, la
% derniere page etant pleine. Fond inchange.
% Lu par main_ensae.tex seul : ne jamais le faire lire par main_institut.tex.
% =====================================================================

\chapter*{Conclusion}
\phantomsection
\addcontentsline{toc}{chapter}{Conclusion}
\markboth{Conclusion}{}

\section*{Apports et résultats}
Ce travail calcule des écarts conditionnels de capital entre états de conformité supposés, puis
examine leur dépendance à la direction de la contagion et au coût des piliers touchés. Sur le
périmètre financier étudié, et non pour un assureur, le besoin de capital lié au risque
opérationnel informatique passe de $6\,049$ à $20\,188$~M\euro{} entre l'état conforme et
l'état non conforme, soit un facteur $3{,}34$. Les deux états diffèrent sur quatre canaux : la
fréquence, la détection, la propagation et l'accumulation. Cet écart, $14\,139$~M\euro{}, a un écart-type de
$734$~M\euro{} entre quatre graines ; il ne mesure pas un effet causal de DORA.
Les canaux se composent presque multiplicativement : l'effet d'une correction dépend de
l'état des autres contrôles. Un plan de remédiation doit donc citer l'effet de fermeture
depuis son état de départ, en incluant les interactions.

La cascade distingue des ordres d'atteinte qu'un score statique des seuls piliers touchés ne représente pas :
$22$ des $26$ ensembles de piliers changent de criticité (note de $1$ à $10$) selon l'ordre d'atteinte.
La normalisation de Leontief garantit la stabilité du système latent, cycles compris. Un processus de Hawkes
ne se distingue pas de la cascade à la résolution disponible, et nous retenons la cascade par parcimonie.

La queue de sévérité et la fréquence d'entrée sont ajustées sur les données ; la propagation et
l'accumulation sont posées. Hors échantillon, la binomiale négative couvre davantage d'années
que le Poisson et la forme de queue n'est pas rejetée, mais douze années donnent peu de puissance.
Le niveau de la charge annuelle est rejeté
(Kolmogorov-Smirnov, $p = 0{,}0042$), sans que la dérive de l'échelle de queue soit confirmée. La
Formule Standard n'exprime pas l'écart entre états : à volume égal, sa charge ne dépend pas de la
conformité, et ses six limites pour cette question (tableau~\ref{tab:critique-formule-standard})
trouvent chacune une réponse conditionnelle dans les chapitres~4 à~10.

% SOURCES-SCRIPTS: 28 39 43 68 92 47 89 91 85 100 74 32

\section*{Non-identifiabilité de la direction}
La direction de la contagion n'est pas identifiable sur les trois sources étudiées
(chronologie américaine, SAS OpRisk, VERIS), pour des raisons distinctes. L'expérience
MOVEit ne distingue pas son effet net de zéro après le placebo. Les deux champs manquants existent
dans le schéma VERIS, mais ne sont renseignés ensemble que dans \emph{un} incident sur $10\,591$.

Le socle sans contagion ne dépend pas de la direction, mais conserve des paramètres posés
(surdispersion, accumulation). Dans le calage de référence (sévérité OpRisk), le capital total se situe entre
$6\,858$ et $8\,697$~M\euro{} sur les $1\,024$ sommets de l'ensemble admissible. Cette plage, large de $1\,839$~M\euro{}, est calculée aux sommets ; recalculée sur quatre graines, sa largeur a un écart-type de $189$~M\euro. Elle n'intègre pas l'incertitude sur la partie symétrique posée par l'expert.

% SOURCES-SCRIPTS: 30 34 28 110 134

\section*{Conséquences pour le régulateur et pour l'entité}
Un registre mieux renseigné permettrait d'éprouver la contagion : survenance à la semaine, classement
par domaine de contrôle, consolidation des causes communes et ensemble des piliers
touchés. Dans le modèle, leur loi jointe peut distinguer une matrice de sa transposée.

La notification DORA pourrait rendre obligatoires les deux premiers champs, déjà
prévus dans VERIS. Fixer les six sens documentés par sept rapports post-incident
réduirait la plage de $32{,}7\,\%$ en moyenne, de $22$ à $44\,\%$ selon la graine,
ordre de grandeur et non promesse. Les dépendances les plus utiles, dont
gouvernance--incidents, en lèveraient chacune $11$ à $13\,\%$, sans ordre résolu.

% SOURCES-SCRIPTS: 66 55 117

Remédier, détenir ou transférer le risque se comparent ensuite en euros par an, avec une prime de réassurance
calculée sur un ratio de sinistralité de $48{,}5\,\%$. Pour l'entité notionnelle de $20\,000$~M\$ d'actifs,
dont le besoin de capital ($24$~M\euro{}) est une extrapolation d'ordre de grandeur, le coût annuel
de détention représente $1{,}35$ à $1{,}71$ fois la sinistralité attendue. Dans le scénario de traité en excès de perte idéalisé, le coût de financement baisse d'environ deux tiers, et de près de moitié pertes conservées comprises ; la capacité du marché, les
exclusions et le risque de contrepartie l'entament. Plusieurs fréquences réelles restent compatibles
avec les mêmes comptes du registre (\S\ref{sec:entites-capture}).



% SOURCES-SCRIPTS: 58 65 67 136

\section*{Limites}
L'élicitation structurée a été préparée et diffusée, mais elle est restée sans réponse (annexe~\ref{ch:elicitation}). La lecture centrale de la direction repose donc sur un barème d'expert
qualitatif. Deux hypothèses pèsent le plus sans être testables sur la donnée : les niveaux de
propagation, posés, et l'additivité des coûts entre piliers touchés. L'élasticité de cette dernière atteint $0{,}95$ à
l'état non conforme contre $0{,}39$ à l'état conforme (\S\ref{sec:additivite-couts}).
Ces sensibilités évaluent l'effet des conventions, sans déterminer leur valeur empirique.

Dans la construction additive, sur les sensibilités étudiées, le signe du surcoût et les deux piliers les plus coûteux, la gouvernance et la gestion des
prestataires tiers, sont conservés. L'ordre de remédiation entre ces deux piliers dépend en revanche du
niveau d'émission de la gouvernance : il tient tant que celle-ci n'est dégradée que du
quart, et bascule à la moitié. Le niveau de capital reste illustratif ; les bandes de calibration
en montrent l'incertitude, conditionnellement aux autres hypothèses.

L'unité de perte modifie l'amplitude et l'attribution de l'écart
(\S\ref{sec:institut-unite-perte}) : une perte unique neutralise financièrement les
canaux agissant sur le nombre de piliers. Le bootstrap conjoint propage l'incertitude
de calibration de la sévérité (\S\ref{sec:institut-bootstrap-joint}) ; ses bandes
restent conditionnelles à la fréquence et aux hypothèses d'expert.

La comparaison de scénarios est l'usage accessible sous ces conventions. Pour chiffrer un
investissement ou arrêter un capital propre à l'entité, il reste à confronter les états et les coûts
par pilier à des observations indépendantes. Un panel renseignerait le barème ; des séquences
d'incidents et des coûts ventilés permettraient d'éprouver la contagion et sa traduction financière.

% SOURCES-SCRIPTS: 74 64 65 119 120
